% SHARD-ID: CA-00-FRONTMATTER
% SHARD-TITLE: Frontmatter, North Star, and Claim Status
% SHARD-SUMMARY: Declares the project north star, the two success criteria, and the evidence rules.
% SHARD-SUMMARY: Fixes the claim-status labels that separate programme-level conjecture from checked results.
% SHARD-KEYWORDS: status, scope, north star, claim status, lab book, evidence chain

\section{Frontmatter, North Star, and Claim Status}
\label{sec:frontmatter}

This is a reproducible research lab book. Its guiding goal --- the
\emph{north star} --- is a single, concrete construction problem:

\begin{quote}
Given a (unitary) fusion or modular tensor category $\Cc$, together with whatever
additional data is needed (OPE coefficients / conformal data), construct a family
of microscopic (lattice) models and their symmetry generators whose continuum
limit \emph{provably} yields a mathematically rigorous QFT/CFT that (i) realises
the input category data and (ii) carries a full projective unitary representation
of the symmetries.
\end{quote}

The working belief is that this is achievable at least for \emph{rational} CFTs,
and perhaps for all QFTs. This is a programme, not a settled theory: every step
toward the goal is a question, proposal, or hypothesis until a local source, a
local derivation, or a reproducible run supports it.

\subsection{The Two Success Criteria}

A construction is only a candidate solution if it meets both, non-negotiably:

\begin{enumerate}
  \item \textbf{Provability.} The continuum limit is established rigorously
  (an operator-algebraic / wavelet renormalisation limit with theorems, not a
  heuristic scaling argument).
  \item \textbf{Fidelity.} The realised continuum content faithfully reproduces
  the input category data $\Cc$, and the symmetry representation is \emph{full}
  and projective-unitary.
\end{enumerate}

\subsection{Claim Status}

No claim is treated as established until it is marked \emph{Source-local} or
\emph{Checked} and linked to a local source, derivation, test, or run artifact.
The status labels are fixed in \texttt{CONVENTIONS.md}: \emph{Question},
\emph{Proposal}, \emph{Source-local}, \emph{Checked}, and \emph{Rejected}. These
labels are part of the scientific record; they separate programme-level
conjecture from results.

\subsection{Evidence Chain}

Every substantive assertion should be traceable through:

\begin{enumerate}
  \item a local source manifest under \texttt{references/} (the kept papers are
  the ground truth),
  \item a convention entry in \texttt{CONVENTIONS.md},
  \item a derivation, script, test, or run bundle when computation is involved,
  \item an index row in \texttt{INDEX.md}, and
  \item a report shard in this lab book.
\end{enumerate}

Missing links are allowed during exploration only when they are explicit.
